\section{Differential lines and section ratios}

We record the elementary differential facts needed to transfer specialness. They apply on ordinary smooth projective manifolds; their root-chart use will be justified explicitly.

\begin{lemma}[Closed-form argument]\label{lem:ratios}
Let $V$ be a connected smooth projective manifold and let $i:L\to\Omega_V^p$ be a nonzero line-bundle map, where $1\le p\le\dim V$. If $s_0,\ldots,s_r$ are sections of the same positive power $L^{\otimes m}$, with $s_0\ne0$, then
\begin{equation}\label{eq:wedge-ratio}
 d(s_a/s_0)\wedge i(L)=0
\end{equation}
at the generic point. In particular, $\kappa(L)\le p$. If $\kappa(L)=p$, there are such sections $s_0,\ldots,s_p$ for which the image line of $i$ is generically spanned by
\begin{equation}\label{eq:wedge-description}
 d(s_1/s_0)\wedge\cdots\wedge d(s_p/s_0)\ne0.
\end{equation}
\end{lemma}

\begin{proof}
Adjoin simultaneous $m$th roots of the finitely many sections in copies of the total space of $L$. The equations $t_a^m=s_a$ are monic and define a finite space over $V$. Choose a reduced irreducible component dominating $V$ and take a smooth projective resolution $q:V'\to V$ of that component \cite{Hironaka}. The tautological roots are holomorphic sections of $q^*L$. Their differential images
\[
 \alpha_a\in H^0(V',\Omega_{V'}^p)
\]
are holomorphic, and $\alpha_0$ is generically nonzero because $q$ is generically finite in characteristic zero.

Holomorphic forms on a compact K\"ahler manifold are closed \cite{Demailly}. On the open set where $t_0\ne0$, we have $\alpha_a=(t_a/t_0)\alpha_0$, so
\[
 d(t_a/t_0)\wedge\alpha_0=0.
\]
Taking the $m$th power of the ratio gives \eqref{eq:wedge-ratio} upstairs; generic injectivity of differential pullback gives it downstairs.

If $r$ independent covectors each wedge to zero with a nonzero $p$-form, expansion in a basis containing those covectors shows that every nonzero term of the form contains all $r$ of them. Hence $r\le p$. In characteristic zero the image dimension of a section map is the generic rank of its ratio differentials, proving the bound. If that dimension is $p$, choose $p$ algebraically independent ratios from one section map. The same exterior-algebra argument gives \eqref{eq:wedge-description}.
\end{proof}

\begin{lemma}[Saturated meromorphic directions]\label{lem:saturation}
Let $E$ be a holomorphic vector bundle on a smooth complex manifold. The intersection with $E$ of a rank-one meromorphic subspace of its meromorphic sections is an invertible sheaf $K$ with a nonzero map $K\to E$. Moreover, a meromorphic tensor in the corresponding line inside $E^{\otimes m}$ is holomorphic if and only if it is a holomorphic section of $K^{\otimes m}$. Both assertions hold equivariantly in smooth root charts.
\end{lemma}

\begin{proof}
In a local frame clear the denominators of a meromorphic generator and divide out all common irreducible factors of its entries. Smooth local rings are factorial, so the resulting vector $v$ has holomorphic entries with no common height-one factor. If $av$ is holomorphic for a meromorphic scalar $a$, a pole of $a$ at a prime divisor would force that prime to divide every entry of $v$. Thus $a$ has no divisorial pole and is holomorphic. This proves that the saturated module is freely generated by $v$.

At every prime divisor some entry of $v$ is a unit at its generic point. The corresponding entry of $v^{\otimes m}$ has the same property. The identical pole test proves the tensor-power assertion. Holomorphic extension across codimension two gives the statements on the whole chart. The construction is intrinsic to the meromorphic direction and hence respects equivariance and changes of chart. The form-valued generator is allowed to vanish in codimension two; the assertion is about an invertible sheaf with a nonzero map, not about a subbundle everywhere.
\end{proof}

\begin{remark}[Comparison of differential conventions]\label{rem:conventions}
Under $z=w^m$, a coefficient $z^{-a}$ in a tensor containing $t$ occurrences of $dz$ pulls back with order
\[
 t(m-1)-ma.
\]
Thus the allowed pole is $a\le\lfloor t(1-1/m)\rfloor$, the usual orbifold symmetric-differential rule. For a fixed generic differential line, saturation of its powers has exactly the invariant sections of the corresponding saturated root line. The primitive-generator argument in Lemma~\ref{lem:saturation} verifies this assertion in codimension one and extends it across codimension two. This compares the section lattices; our specialness argument uses the line-bundle convention of Definition~\ref{def:special} directly.
\end{remark}

\begin{lemma}[Rational descent]\label{lem:rational-descent}
Let $g:Y\to X$ be a dominant morphism of connected smooth projective varieties with connected general fibre. A rational function on $Y$ which is constant on general fibres, wherever defined, is the pullback of a rational function on $X$.
\end{lemma}

\begin{proof}
Let $Z\subset X\times\PP^1$ be the closure of the joint image of $g$ and the function, on its regular locus. The actual image is constructible and dense in $Z$, so it contains a dense open of $Z$. Its fibres over general points of $X$ have at most one point. The fibre-dimension theorem therefore makes $Z\to X$ generically finite. After shrinking $X$, its generic degree is realized by distinct points, since the characteristic is zero, and all those points belong to the actual-image open. The degree is consequently one. Hence $\CC(Z)=\CC(X)$, giving the descent. Connectedness of the general fibre ensures that the constant is a single value on that fibre.
\end{proof}
